\section{4.8.Gestion de errores}\label{gestion-de-errores}

\subsection{4.8. Gestión de errores y robustez de la
solución}\label{gestiuxf3n-de-errores-y-robustez-de-la-soluciuxf3n}

Debido a la dependencia de herramientas externas y a la complejidad del
proceso de análisis, durante el diseño de EVMAudit se prestó especial
atención a la gestión de errores y a la robustez del sistema.

La existencia de múltiples componentes externos hace necesario disponer
de mecanismos que permitan detectar y manejar situaciones anómalas sin
comprometer el funcionamiento global de la herramienta.

\subsubsection{4.8.1. Necesidad de una gestión
centralizada}\label{necesidad-de-una-gestiuxf3n-centralizada}

Durante el proceso de análisis pueden producirse diferentes tipos de
errores:

\begin{itemize}
\tightlist
\item
  ausencia de herramientas instaladas
\item
  incompatibilidades entre versiones
\item
  contratos inválidos
\item
  timeouts durante el análisis
\item
  errores internos de las herramientas utilizadas
\item
  problemas durante la generación de wrappers.
\end{itemize}

La gestión individual de cada una de estas situaciones complicaría
considerablemente el mantenimiento del sistema.

Por este motivo se implementó una jerarquía de excepciones propia que
centraliza el tratamiento de los errores.

\subsubsection{4.8.2. Jerarquía de
excepciones}\label{jerarquuxeda-de-excepciones}

La arquitectura implementada incorpora una clase base común a partir de
la cual se derivan los distintos tipos de error utilizados por EVMAudit.

Esta aproximación permite distinguir claramente entre:

\begin{itemize}
\tightlist
\item
  errores de configuración
\item
  errores de análisis
\item
  errores producidos por herramientas externas.
\end{itemize}

La utilización de excepciones específicas simplifica la depuración y
mejora la legibilidad del código.

!TODO: insertar diagrama UML de excepciones.

\subsubsection{4.8.3. Detección de dependencias
externas}\label{detecciuxf3n-de-dependencias-externas}

Antes de ejecutar las herramientas integradas, EVMAudit verifica que las
dependencias necesarias se encuentren correctamente instaladas.

Este mecanismo permite detectar problemas de configuración de forma
temprana y proporcionar mensajes de error más descriptivos.

La validación previa evita que el proceso de análisis falle de forma
inesperada en fases posteriores.

\subsubsection{4.8.4. Gestión de timeouts}\label{gestiuxf3n-de-timeouts}

Algunas herramientas, especialmente aquellas basadas en ejecución
simbólica, pueden presentar tiempos de análisis elevados dependiendo de
la complejidad del contrato.

Con el objetivo de evitar bloqueos indefinidos, EVMAudit incorpora
límites temporales para determinadas operaciones.

La utilización de timeouts permite:

\begin{itemize}
\tightlist
\item
  mejorar la estabilidad del sistema
\item
  evitar consumos excesivos de recursos
\item
  garantizar la finalización del análisis.
\end{itemize}

\subsubsection{4.8.5. Preservación de resultados ante
errores}\label{preservaciuxf3n-de-resultados-ante-errores}

Cuando se produce un fallo durante alguna fase del análisis, EVMAudit
intenta conservar la información obtenida hasta ese momento.

Esta estrategia permite:

\begin{itemize}
\tightlist
\item
  facilitar la depuración
\item
  conservar evidencias parciales
\item
  evitar la pérdida completa del trabajo realizado.
\end{itemize}

La preservación de resultados resulta especialmente relevante en
entornos de auditoría, donde incluso los análisis incompletos pueden
aportar información útil.

\subsubsection{4.8.6. Robustez frente a comportamientos específicos de
las
herramientas}\label{robustez-frente-a-comportamientos-especuxedficos-de-las-herramientas}

Durante el desarrollo se detectaron determinados comportamientos
particulares en las herramientas integradas.

Entre ellos destacan:

\begin{itemize}
\tightlist
\item
  códigos de retorno no convencionales
\item
  formatos de salida inconsistentes
\item
  diferencias entre versiones
\item
  mensajes mezclados con datos estructurados.
\end{itemize}

Para hacer frente a estas situaciones se implementaron mecanismos
específicos de adaptación y validación.

Estas medidas permitieron integrar las herramientas seleccionadas sin
modificar su funcionamiento interno.

\subsubsection{4.8.7. Beneficios de la estrategia
adoptada}\label{beneficios-de-la-estrategia-adoptada}

La incorporación de mecanismos de gestión de errores proporciona
diversas ventajas:

\begin{itemize}
\tightlist
\item
  mayor estabilidad
\item
  mejora de la experiencia de usuario
\item
  facilidad de mantenimiento
\item
  simplificación de la depuración
\item
  incremento de la robustez general del sistema.
\end{itemize}

La existencia de una capa de gestión de errores independiente contribuye
además a mantener desacoplados los distintos módulos que forman la
arquitectura de EVMAudit.

!TODO: insertar referencia cruzada al módulo Runner.

!TODO: insertar referencia cruzada al diagrama de arquitectura general.
